\section{Lecture 15 (February 19, 2025, est.): Schur functors and Deligne's theorem}
\label{lec15}
\noindent\textit{Scribe: Guillermo Sanmarco}

Throughout this lecture, $k$ is algebraically closed of characteristic
zero. We follow the framework of \cite{Deligne2002}. For part of the
discussion, one may drop the finite-length and finite-dimensional-Hom
conditions: retain an essentially small $k$-linear abelian rigid symmetric
monoidal category $\mathcal C$ with $\End(\1)=k$. Tensor product is exact
in each variable, because rigidity supplies both adjoints. These are
still substantive hypotheses; retaining rigidity alone would not suffice.

\subsection{The statement and the finiteness condition}

\begin{definition}
An object $X\in\mathcal C$ is of \emph{moderate growth} if its tensor
powers have finite length and there is $a_X\geq1$ with
\[
 \operatorname{length}(X^{\otimes n})\leq a_X^n
       \qquad(n\geq1).
\]
It is \emph{Schur-finite} if $S^\lambda(X)=0$ for some partition
$\lambda$, where the Schur functor is defined below.
\end{definition}

\begin{theorem}[Deligne]
An object is Schur-finite if and only if it is of moderate growth.
Schur-finite objects are closed under finite direct sums, tensor products,
duals, extensions and subquotients. A symmetric tensor category has a
fiber functor to $\SVec_k$ if and only if all its objects are Schur-finite.
\end{theorem}

The existence of the super fiber functor is the deep part of this
theorem. This lecture establishes the Schur decomposition and the
implication from moderate growth to Schur-finiteness, and records the
representation-theoretic tools used in the full proof. The remaining
arguments are in \cite[Proposition 0.5 and Theorem 0.6]{Deligne2002};
they are not consequences of the growth estimate alone.

\begin{lemma}
If every object of $\mathcal C$ has finite length, then every morphism
space is finite-dimensional.
\end{lemma}
\begin{proof}
We use the standard simple-unit lemma for rigid abelian tensor
categories: every subobject of $\1$ is a direct summand, so
$\End(\1)=k$ implies that $\1$ is simple
\cite[Proposition 1.17]{DMTannakian}. This is the input used in
\cite[Proposition 1.1]{Deligne2002}; we do not prove that lemma here.
For an object $Z$ of finite length, a linearly independent family of
maps $\1\to Z$ gives a monomorphism from the corresponding finite
direct sum of copies of $\1$ to $Z$: any kernel is a direct sum of
copies of the simple object $\1$ and would give a linear relation.
Thus $\dim_k\Hom(\1,Z)\leq\operatorname{length}(Z)$.
Finally rigidity gives
\[
 \Hom(X,Y)\cong\Hom(\1,X^*\otimes Y),
\]
and $X^*\otimes Y$ has finite length by assumption.
\end{proof}

\subsection{Finite-group actions inside a category}

For a finite-dimensional vector space $V$ and $X\in\mathcal C$, the
notation $V\otimes X$ denotes the corresponding finite direct sum of
copies of $X$, characterized independently of a choice of basis.
Define
\[
 \underline{\Hom}(V,X)=V^*\otimes X.
\]
Suppose a finite group $G$ acts on $X$, that is, we have a homomorphism
$G\to\Aut_{\mathcal C}(X)$. The averaging operator
\[
 e_G=\frac1{|G|}\sum_{g\in G}g\in\End_{\mathcal C}(X)
\]
is idempotent. Its image is the invariant object $X^G$. Its kernel is
the sum of the images of $g-1$, so the coinvariant object
$X_G=\operatorname{coker}(\bigoplus_{g\in G}X\xrightarrow{(g-1)}X)$
is canonically isomorphic to $\im e_G$.

If $V\in\Rep_k(G)$, there is a diagonal action on $V^*\otimes X$,
and we put
\[
 \underline{\Hom}_G(V,X)=(V^*\otimes X)^G.
\]
This is an object of $\mathcal C$, rather than merely the vector space
of morphisms from $V$ to $X$.

\begin{proposition}[Isotypic decomposition]
Let $(V_\lambda)_{\lambda\in I}$ be representatives of the irreducible
$kG$-modules. For every object $X$ carrying a $G$-action, evaluation
induces a natural $G$-equivariant isomorphism
\[
 \bigoplus_{\lambda\in I}
 V_\lambda\otimes\underline{\Hom}_G(V_\lambda,X)
       \xrightarrow{\sim}X.
\]
\end{proposition}
\begin{proof}
Maschke's theorem and the algebraic closure of $k$ give
$kG\cong\prod_{\lambda\in I}\End_k(V_\lambda)$.
The central idempotents split $X$ into isotypic summands. In the
$\lambda$-summand, matrix units split it into $\dim V_\lambda$
isomorphic summands; their multiplicity object is precisely
$(V_\lambda^*\otimes X)^G$. The resulting isomorphism is evaluation,
so it is independent of the matrix units chosen.
\end{proof}

\subsection{Partitions and Schur functors}

A partition of $n$ is a finite sequence
$\lambda=(\lambda_1,\ldots,\lambda_r)$ with
$\lambda_1\geq\cdots\geq\lambda_r>0$ and
$|\lambda|=\sum_i\lambda_i=n$. Its Young diagram is
\[
 [\lambda]=\{(i,j):1\leq i\leq r,\ 1\leq j\leq\lambda_i\}.
\]
The transpose $\lambda^{\mathsf t}$ is defined by interchanging the
two coordinates. For example, $(3,1)^{\mathsf t}=(2,1,1)$.
We also allow the empty partition of $0$, with Schur functor equal
to the constant unit object; this convention is used in the direct-sum
formula below.
The irreducible representations of $S_n$ are indexed by partitions of
$n$; write $V_\lambda$ for the corresponding representation.

The symmetry on $\mathcal C$ defines an action of $S_n$ on
$X^{\otimes n}$ by adjacent transpositions. The symmetric, rather
than merely braided, condition ensures that these transpositions
square to the identity.

\begin{definition}
For $\lambda\vdash n$, the \emph{Schur functor} is
\[
 S^\lambda(X)=
 \underline{\Hom}_{S_n}(V_\lambda,X^{\otimes n}).
\]
\end{definition}

Applying the isotypic decomposition gives
\begin{equation}\label{eq:lec15-schur-decomposition}
 X^{\otimes n}\cong
 \bigoplus_{\lambda\vdash n}V_\lambda\otimes S^\lambda(X).
\end{equation}
For the one-row and one-column partitions,
$V_{(n)}$ is trivial and $V_{(1^n)}$ is the sign representation.
Consequently
\[
 S^{(n)}(X)=\Sym^n(X),\qquad
 S^{(1^n)}(X)=\bigwedge^nX.
\]
In characteristic zero these are the images of the symmetrizing and
antisymmetrizing idempotents, respectively.

\begin{proposition}
If $X$ is of moderate growth, then some Schur functor annihilates $X$.
\end{proposition}
\begin{proof}
Suppose instead that $S^\lambda(X)\ne0$ for every partition $\lambda$.
Each summand in \eqref{eq:lec15-schur-decomposition} then contributes
at least $\dim V_\lambda$ to the length. Therefore
\[
 \operatorname{length}(X^{\otimes n})
 \geq\sum_{\lambda\vdash n}\dim V_\lambda
 \geq\left(\sum_{\lambda\vdash n}(\dim V_\lambda)^2\right)^{1/2}
 =\sqrt{n!}.
\]
The last equality is the dimension of the regular representation of
$S_n$. For even $n$, $n!\geq(n/2)^{n/2}$, so
$(\sqrt{n!})^{1/n}\geq(n/2)^{1/4}$ tends to infinity.
This contradicts any exponential bound $a_X^n$.
\end{proof}

\begin{example}
For an ordinary vector space of dimension $d$, $\bigwedge^{d+1}V=0$.
An odd one-dimensional supervector space instead has
$\Sym^2(V)=0$, while its exterior powers never vanish. Thus exterior
powers alone cannot characterize super fiber functors. Schur functors
allow both kinds of vanishing.
\end{example}

\subsection{Restriction, induction, and multiplicities}

For $n=n_1+\cdots+n_r$, the Young subgroup
$S_{n_1}\times\cdots\times S_{n_r}\subseteq S_n$ acts on consecutive
blocks of letters. Its irreducibles are the external tensor products
$V_{\mu_1}\boxtimes\cdots\boxtimes V_{\mu_r}$, with
$\mu_i\vdash n_i$. Define
\[
 [\lambda:\mu_1,\ldots,\mu_r]
 =\dim_k\Hom_{\prod_iS_{n_i}}
 \left(\boxtimes_iV_{\mu_i},
       \operatorname{Res}^{S_n}_{\prod_iS_{n_i}}V_\lambda\right).
\]
Frobenius reciprocity identifies this with the multiplicity of
$V_\lambda$ in
$\operatorname{Ind}^{S_n}_{\prod_iS_{n_i}}(\boxtimes_iV_{\mu_i})$.
These nonnegative integers are computed by the
Littlewood--Richardson rule, iterated when $r>2$.

For example, when $r=2$ the same decomposition yields
\[
 S^\lambda(X\oplus Y)\cong
 \bigoplus_{|\mu|+|\nu|=|\lambda|}
 \left(S^\mu(X)\otimes S^\nu(Y)\right)^%
 {\oplus[\lambda:\mu,\nu]}.
\]
Indeed, expand $(X\oplus Y)^{\otimes n}$ by the subsets of tensor
positions occupied by $X$; each fixed cardinality is induced from its
Young subgroup, and then apply the multiplicity-object functor.
Likewise,
\[
 S^\mu(X)\otimes S^\nu(X)\cong
 \bigoplus_{\lambda\vdash |\mu|+|\nu|}
 S^\lambda(X)^{\oplus[\lambda:\mu,\nu]}.
\]
These formulas explain how the representation theory of symmetric
groups controls the operations appearing in the Schur-finiteness
criterion.
